% Generated draft: requires mathematical and source review.
\documentclass[11pt]{article}
\usepackage[margin=25mm]{geometry}
\usepackage{fontspec}
\setmainfont{Latin Modern Roman}
\usepackage{amsmath,amssymb,mathtools,unicode-math}
\setmathfont{Latin Modern Math}
\usepackage{enumitem,needspace,xurl,hyperref}
\setlength{\parindent}{0pt}
\setlength{\parskip}{5pt}
\setlength{\emergencystretch}{3em}
\newcommand{\N}{\mathbb{N}}
\newcommand{\Z}{\mathbb{Z}}
\newcommand{\Q}{\mathbb{Q}}
\newcommand{\R}{\mathbb{R}}
\newcommand{\C}{\mathbb{C}}
\newcommand{\F}{\mathbb{F}}
\newcommand{\A}{\mathbb{A}}
\newcommand{\PP}{\mathbb{P}}
\newcommand{\E}{\mathbb{E}}
\newcommand{\eps}{\varepsilon}
\newcommand{\dd}{\,\mathrm{d}}
\begin{document}
\section*{UnsolvedMath ID: 2000013}
\textbf{Source problem code:} catalogue entry
\subsection{Definitions and mathematical statement}
A dynamical description of an elastic evolution uses a phase space $\mathcal X$ of admissible states, conventionally deformation and velocity pairs $(y,v)$. When a chosen autonomous evolution is globally well posed, its solution maps $S(t):\mathcal X\to\mathcal X$ satisfy
\[
S(0)=\operatorname{id},\qquad S(t+s)=S(t)\circ S(s)\quad(t,s\geq0).
\]
Qualitative dynamics concerns structural properties of these trajectories, such as invariant sets, stability, and long-time behavior, rather than only construction of individual solutions. If solutions are not unique or not global, a different evolution framework may be necessary; no semiflow is assumed to exist for general dynamic elasticity.

\textbf{Research problem.} Develop a qualitative dynamics for dynamic theories of elasticity. In a specified theory, this requires selecting the evolution equations, admissible initial and boundary data, solution notion, and phase-space topology before formulating a precise dynamical assertion. The supplied record specifies none of these. Its research summary gives global attractors and long-time classification as examples of the desired analysis, and notes partial tools in damped or viscoelastic settings. Those examples are not an assertion that every conservative or nondissipative elastic system must possess an attractor.

For reference, a global attractor for a given continuous semiflow is a compact invariant set $\mathcal A$ that attracts every bounded set in the chosen phase-space topology. This is a definition of one possible target, not an added requirement that the source problem be solved by proving such an object exists. The summary describes the broad general program as open, with moderate-confidence OPEN-TRIAGE status.
\subsection{Short English statement}
Build a framework for stability and long-time behavior in dynamic elasticity. The supplied request is a research program rather than a theorem with specified equations, phase space, or hypotheses.
\subsection{Sources}
\begin{itemize}
\item J. M. Ball, \emph{Some Open Problems in Elasticity} (2002), Problem 13, \url{https://people.maths.ox.ac.uk/ball/Articles%20in%20Conference%20Proceedings%20and%20Books/JMB%202002%20re%20Marsden%2060th.pdf}; PDF-text extraction, accessed 2026-07-29.
\item The supplied research summary mentions partial tools for damping and viscoelasticity but gives no specific publication or theorem citation for those tools or its current-status assessment.
\item Local registry \texttt{data/unsolvedmath/problems.json}, permanent ID $2000013$, source code catalogue entry, \emph{Some Open Problems in Elasticity --- Qualitative dynamics}. The record labels this open and OPEN-TRIAGE, with NEEDS\_REVIEW source-status and rights qualifications.
\end{itemize}
\end{document}
